\section{Camera branch}
The camera branch tests two representations of visual observations while
keeping the Phase~1 windows and R1/R2 production code unchanged. Both easy
Machine Hall flights provide cam0 grayscale images at approximately
$20\,\mathrm{Hz}$. Verification opened all 3,682 MH\_01 and 3,040 MH\_02
images and found 40 frames inside every two-second window. No missing image
or uncovered window was accepted.

\subsection{Offline preprocessing and Tier 1}
Images are converted once to $64\times64$ uint8 arrays. Integer nanosecond
timestamps select the first, nearest-midpoint and last frames strictly
inside the half-open window interval. The resulting tensors have shape
$(720,3,64,64)$ and $(592,3,64,64)$, in the original window CSV order.
Only raw images remain outside version control; small arrays, indices and
checksum manifests are supplied for the other branches. R1's existing
dataset divides these arrays by 255 to produce float32 inputs.

Tier~1 uses Farneb\"ack dense forward flow between consecutive resized
frames~\cite{r3farneback2003}. Displacements are divided by the actual frame
interval to obtain resized-image pixels per second. A window aggregates
only pairs whose two frames lie inside it. Ten predictors describe mean,
standard deviation and mean per-pair 90th-percentile flow magnitude; signed
horizontal/vertical mean and standard deviation; mean and standard
deviation of Laplacian variance; and mean Shi--Tomasi corner
count~\cite{r3shi1994}. IDs, timestamps and frame counts are excluded from
the predictors. Feature extraction does not read labels, ground truth,
physics or IMU data. Flow is an image-space observation, not a metric
velocity measurement; sharpness and corners also depend on scene texture.

The existing Phase~1 Lasso tuning code is reused unchanged, with its
five purged time blocks, alpha grid, median-MAE selection and in-pipeline
standardisation. Each variant is tuned on the training flight, refitted on
that whole flight and evaluated on the other flight. Table~\ref{tab:r3tier1}
reports the actual results. The 104-feature baseline reproduces mean MAE
$2.50\,\mathrm{cm}$. Appending ten camera predictors increases the mean to
$3.63\,\mathrm{cm}$. On the MH\_02 test fold, training-only CV selects
$\alpha=0.1$, shrinking all coefficients to zero; this is a failure of this
representation/tuning combination to improve cross-flight performance, not
evidence that all visual representations are uninformative.

\begin{table}[t]
\centering
\caption{Tier 1 held-out-flight results. MAE is in centimetres.}
\label{tab:r3tier1}
\begin{tabular}{llrr}
\hline
Inputs & Test flight & MAE & $R^2$ \\
\hline
104 IMU features & MH\_01 & 2.59 & 0.676 \\
104 + 10 camera & MH\_01 & 2.62 & 0.673 \\
104 IMU features & MH\_02 & 2.42 & 0.837 \\
104 + 10 camera & MH\_02 & 4.65 & $-0.022$ \\
\hline
\end{tabular}
\end{table}

\subsection{Tier 2 architecture and evaluation}
FlowNet motivates learning from temporally related images, but its
supervised flow network and training corpus are not used
here~\cite{r3dosovitskiy2015}. E5 stacks start/middle/end grayscale frames as
three temporal channels. Three $3\times3$, stride-2 convolution blocks with
16, 32 and 64 channels each apply four-group GroupNorm, ReLU and spatial
dropout. GroupNorm uses per-example statistics instead of batch running
statistics~\cite{r3wu2018}. Global average pooling produces the shared
64-number embedding. A pre-pooling activation has a $15\times15$ receptive
field. The encoder has 23,808 parameters and R2's unchanged linear head
adds 65 for magnitude prediction. A four-output head is supported by the
architecture but is not trained in E5.

E5 uses R1's fixed partitions: 450/119 training/validation windows when
testing MH\_01 and 553/144 when testing MH\_02, with the four-second temporal
gap. AdamW, MSE, learning rate $10^{-3}$, weight decay $10^{-4}$, dropout 0.3,
the validation plateau schedule and early stopping follow the existing DL
configuration. The epoch limit is 100, stopping patience is 15 and minimum
improvement is $10^{-5}$. R2's current trainer hard-codes IMU inputs, so an
R3-specific runner reuses its seed/result helpers without modifying that
trainer. Test tensors are evaluated after the best validation weights are
restored. Seeds 42, 43 and 44 are all reported; none is selected by its test
score. Restored epochs are 11/11/18 and 24/40/17 for the two test flights.

\begin{table}[t]
\centering
\caption{Camera-only E5: mean $\pm$ sample standard deviation over three
seeds. MAE and RMSE are in centimetres.}
\label{tab:r3e5}
\begin{tabular}{lrrr}
\hline
Test flight & MAE & RMSE & $R^2$ \\
\hline
MH\_01 & $3.79\pm0.11$ & $5.46\pm0.22$ & $0.100\pm0.073$ \\
MH\_02 & $5.05\pm0.21$ & $7.67\pm0.12$ & $0.047\pm0.031$ \\
\hline
\end{tabular}
\end{table}

Equal-fold MAE, computed per seed and then averaged, is
$4.42\pm0.12\,\mathrm{cm}$. This is close to the Phase~1 mean baseline
($4.53\,\mathrm{cm}$) and worse than Lasso ($2.50\,\mathrm{cm}$). The
$R^2$ values indicate weak explanatory power. Validation MAE at the selected
epochs is 1.72--1.95 cm and 2.67--2.85 cm, below the test-flight errors,
showing limited transfer. Training curves remain finite and early stopping
restores the validation-selected weights; regularisation does not guarantee
cross-flight accuracy. Training losses include dropout, so their magnitude
is not directly comparable to evaluation-mode validation losses.

The measured conclusion is that these Tier~1 features do not improve the
Lasso baseline and this E5 image-only model does not outperform it. Whether
the 64-dimensional camera embedding adds complementary information to IMU
embeddings remains E6/A3, owned by the fusion/evaluation branches. Three
seeds quantify optimisation spread, not uncertainty over a wider flight
population. Lasso refits on the full training flight while E5 reserves a
validation block; this protocol difference must accompany comparisons.
Finally, the Phase~1 caveat remains: labels and physics priors initialise
windows from ground-truth state, although camera preprocessing never uses it.
